% !TeX root = ../main.tex
\section{Opis oprogramowania}
\label{sec:radar_implementation}

% ------------------------------------------------------------------------
% Rozdział 3: Opis oprogramowania
% 3.1 Środowisko symulacyjne MATLAB i wykorzystane pakiety narzędziowe
% 3.2 Ogólna architektura oprogramowania i przepływ danych
% 3.3 Implementacja nadajnika sygnału IEEE 802.11
% 3.4 Implementacja modelu kanału propagacyjnego i celów ruchomych
% 3.5 Implementacja odbiornika radarowego (3.5.1-3.5.4)
% ------------------------------------------------------------------------

\subsection{Środowisko symulacyjne MATLAB i~wykorzystane pakiety narzędziowe}
\label{subsec:matlab_environment_toolboxes}

Symulator pasywnego radaru wykorzystującego sygnały IEEE~802.11 zaimplementowano w~środowisku obliczeniowym MATLAB (weryfikowanym w~wydaniach R2024b oraz R2025b). Wybór tego środowiska podyktowano jego natywną optymalizacją pod kątem operacji wektorowych i~macierzowych na liczbach zespolonych podwójnej precyzji, co bezpośrednio odpowiada strukturze danych w~przetwarzaniu sygnałów OFDM. Podstawowymi obiektami obliczeniowymi w~torze odbiorczym są dwuwymiarowe macierze podnośnych o~wymiarach $N_{\mathrm{FFT}} \times M$ oraz macierze periodogramów o~wymiarach $N_{\mathrm{per}} \times M_{\mathrm{per}}$, na których wykonywane są szybkie transformaty Fouriera (\texttt{fft}, \texttt{ifft}) z~dopełnianiem zerami, operacje przesunięcia widma (\texttt{fftshift}, \texttt{ifftshift}) oraz indeksowanie logiczne. Takie podejście pozwoliło na zwięzłą, wektoryzowaną implementację bez konieczności ręcznego zarządzania pamięcią czy pisania zagnieżdżonych pętli w~językach niższego poziomu.

Konfigurację ścieżek uruchomieniowych zautomatyzowano w~skrypcie \texttt{init\_paths.m}, umieszczonym w~katalogu głównym repozytorium. Skrypt ten dynamicznie wyznacza bezwzględną ścieżkę bazową projektu (wykorzystując funkcję \texttt{mfilename}) i~dodaje do ścieżki wyszukiwania MATLAB-a katalogi modułów potoku (\texttt{src/}), biblioteki funkcji numerycznych (\texttt{lib/}) oraz rekurencyjnie podkatalogi skryptów eksperymentalnych i~wizualizacyjnych (\texttt{scripts/}). W~celu zagwarantowania pełnej powtarzalności eksperymentów symulacyjnych, wykresów oraz automatycznych testów regresyjnych wprowadzono deterministyczną inicjalizację generatora liczb pseudolosowych za pomocą polecenia \texttt{rng(seed, 'twister')} (algorytm Mersenne Twister). Pozwoliło to na wierne odtwarzanie losowych ciągów bitów pola danych, realizacji statystycznego kanału wielodrogowego oraz próbek szumu termicznego między kolejnymi uruchomieniami.

Istotnym założeniem projektowym był ścisły rozdział między wykorzystaniem gotowych bibliotek firmy MathWorks a~autorską implementacją algorytmów cyfrowego przetwarzania sygnałów. Z~pakietów narzędziowych (ang.~\textit{Toolboxes}) skorzystano wyłącznie w~zakresie, w~jakim dostarczają one zweryfikowanych modeli standardów telekomunikacyjnych lub podstawowych prymitywów sygnałowych:
\begin{itemize}
    \item \textbf{WLAN Toolbox} -- użyty do zgodnej ze standardami IEEE 802.11a i~802.11ax syntezy przebiegów w~paśmie podstawowym (obiekty konfiguracyjne \texttt{wlanNonHTConfig} oraz \texttt{wlanHESUConfig}, funkcje metadanych \texttt{wlanHEOFDMInfo}, \texttt{wlanFieldIndices}, \texttt{wlanSampleRate} i~\texttt{transmitTime}, a~także generator \texttt{wlanWaveform\-Generator}), standaryzowanego modelowania wielodrogowego kanału wewnątrzbudynkowego (\texttt{wlanTGaxChannel}) oraz wstępnej synchronizacji czasowej pakietu w~odbiorniku (\texttt{wlanPacketDetect}, \texttt{wlanSymbol\-TimingEstimate}).
    \item \textbf{Communications Toolbox} -- wykorzystany w~module kanałowym (\texttt{src/\allowbreak step2\_channel.m}) do kontrolowanego dodawania zespolonego addytywnego szumu białego Gaussa funkcją \texttt{awgn}, z~jawnym zakotwiczeniem poziomu szumu względem mocy nominalnej sygnału nadawczego.
    \item \textbf{Signal Processing Toolbox} -- zastosowany wyłącznie do wygenerowania jednowymiarowych wektorów wag czterowyrazowego okna Blackmana--Harrisa za pomocą funkcji \texttt{blackmanharris} w~pliku \texttt{lib/\allowbreak compute\_\allowbreak range\_\allowbreak doppler\_\allowbreak map.m}. Normalizacja energetyczna wag oraz złożenie separowalnego okna dwuwymiarowego ograniczonego do pasma aktywnych podnośnych zostały zrealizowane jako kod własny.
\end{itemize}

W~projekcie świadomie zrezygnowano z~wyspecjalizowanych pakietów radiolokacyjnych, takich jak \textit{Radar Toolbox} czy \textit{Phased Array System Toolbox}. Cały tor przetwarzania radarowego -- obejmujący rafinację synchronizacji metodą detekcji pierwszego czoła, demodulację OFDM po usunięciu prefiksu cyklicznego, estymację częstotliwościowej odpowiedzi kanału metodą wymuszania zera, filtrację ech stałych w~czasie wolnym, adaptacyjną interpolację luki DC, wyznaczanie mapy odległość--prędkość wraz z~kalibracją osi fizycznych, progowanie CFAR oraz iteracyjną ekstrakcję celów algorytmami Binary~STC i~Coherent~STC -- został napisany od podstaw w~czystym języku MATLAB. Uniknięcie gotowych bloków typu ,,czarna skrzynka'' zapewniło pełną przejrzystość matematyczną implementacji, bezpośrednią zgodność z~modelem teoretycznym opisanym w~rozdziale~\ref{sec:passive_radar} oraz dostęp do wszystkich macierzy pośrednich na potrzeby analizy diagnostycznej i~eksperymentów z~rozdziału~\ref{sec:simulation_results}. Zestawienie podziału kompetencji pomiędzy pakiety narzędziowe a~moduły własne przedstawiono w~tabeli~\ref{tab:toolboxes_vs_custom}.

\begin{table}[H]
    \centering
    \linespread{1.05}\selectfont
    \caption{Podział funkcjonalności symulatora pomiędzy pakiety narzędziowe środowiska MATLAB a~autorską implementację DSP}
    \label{tab:toolboxes_vs_custom}
    \footnotesize
    \begin{tabular}{|p{2.6cm}|p{4.8cm}|p{2.1cm}|p{4.9cm}|}
        \hline
        \textbf{Warstwa / Pakiet} & \textbf{Wykorzystane funkcje / Moduły} & \textbf{Etap potoku} & \textbf{Rola w~symulatorze} \\
        \hline
        MATLAB (rdzeń) & \texttt{fft}, \texttt{ifft}, \texttt{fftshift}, \texttt{ifftshift}, \texttt{rng}, \texttt{init\_paths.m} & Cały potok & Operacje macierzowe na liczbach zespolonych, transformaty 1D/2D, zarządzanie ścieżkami i~determinizm obliczeń \\
        \hline
        WLAN Toolbox & \texttt{wlanNonHTConfig}, \texttt{wlanHESUConfig}, \texttt{wlanHEOFDMInfo}, \texttt{wlanFieldIndices}, \texttt{wlanSampleRate}, \texttt{transmitTime}, \texttt{wlanWaveform\-Generator} & Nadajnik (\texttt{step1}) & Zgodna ze standardami IEEE~802.11a/ax parametryzacja i~generacja ramek PPDU w~paśmie podstawowym \\
        \hline
        WLAN Toolbox & \texttt{wlanTGaxChannel} & Kanał (\texttt{step2}) & Symulacja wielodrogowego kanału wewnątrzbudynkowego (profile Model-A/B/C/E) jako źródła DPI i~ech stałych \\
        \hline
        WLAN Toolbox & \texttt{wlanPacketDetect}, \texttt{wlanSymbol\-TimingEstimate} & Odbiornik (\texttt{step3}) & Zgrubne wykrycie pakietu (L-STF) oraz wyznaczenie metryki korelacyjnej dla preambuły (L-LTF) \\
        \hline
        Communications Toolbox & \texttt{awgn} & Kanał (\texttt{step2}) & Dodawanie zespolonego szumu termicznego AWGN o~zadanym stosunku SNR względem mocy odniesienia \\
        \hline
        Signal Processing Toolbox & \texttt{blackmanharris} & Odbiornik (\texttt{step3}) & Wyznaczenie jednowymiarowych wag okna Blackmana--Harrisa przed autorską normalizacją i~złożeniem okna 2D \\
        \hline
        Implementacja własna (\texttt{src/}, \texttt{lib/}) & \texttt{get\_wifi\_params}, \texttt{configure\_targets}, \texttt{demodulate}, \texttt{estimate\_channel}, \texttt{interpolate\_\allowbreak dc\_nulls}, \texttt{compute\_\allowbreak range\_\allowbreak doppler\_map}, \texttt{calculate\_\allowbreak physical\_axes}, \texttt{binary\_*}, \texttt{coherent\_*}, \texttt{compute\_pnfr} & Cały tor radarowy (\texttt{step1}--\texttt{step4}) & Skalowanie ramek, synteza celów ruchomych z~opóźnieniem ułamkowym, synchronizacja \textit{leading-edge}, demodulacja, ZF, MTI, interpolacja DC, periodogram 2D, CFAR, Binary/Coherent~STC i~metryka PNFR \\
        \hline
    \end{tabular}
\end{table}

\subsection{Ogólna architektura oprogramowania i~przepływ danych}
\label{subsec:software_architecture}

Oprogramowanie symulatora zorganizowano w~architekturze dwupoziomowej. Poziom nadrzędny stanowią cztery funkcje orkiestrujące z~katalogu \texttt{src/}, odwzorowujące fizyczny podział systemu radarowego na nadajnik (\texttt{step1\_\allowbreak transmitter}), kanał propagacyjny (\texttt{step2\_\allowbreak channel}) oraz odbiornik, przy czym odbiornik rozdzielono na część przetwarzania sygnału (\texttt{step3\_\allowbreak receiver}) i~część detekcji celów (\texttt{step4\_\allowbreak interpreter}). Funkcje te pełnią rolę fasad (ang.~\textit{facade}): odpowiadają za rozwiązanie argumentów opcjonalnych i~wartości domyślnych, kolejność wywołań oraz złożenie struktur wyjściowych, nie zawierając jednak samych obliczeń numerycznych. Poziom podrzędny tworzy biblioteka \texttt{lib/}, w~której każdą elementarną operację -- parametryzację ramki, demodulację OFDM, estymację transmitancji kanału, interpolację luki DC, wyznaczenie periodogramu, kalibrację osi czy detekcję celów -- wydzielono do osobnej funkcji o~jawnie zdefiniowanych wejściach i~wyjściach, bez zmiennych globalnych i~bez efektów ubocznych poza komunikatami diagnostycznymi. Taki podział umożliwił niezależne testowanie poszczególnych bloków oraz ich wielokrotne wykorzystanie w~różnych konfiguracjach eksperymentów.

Nad potokiem głównym znajduje się warstwa sterująca, obejmująca skrypt konfiguracji ścieżek \texttt{init\_paths.m}, skrypt demonstracyjny \texttt{scripts/\allowbreak run\_random\_demo.m}, generatory rysunków pracy (\texttt{scripts/\allowbreak figures/}), skrypty badań statystycznych (\texttt{scripts/\allowbreak experiments/\allowbreak run\_\allowbreak detection\_\allowbreak monte\_\allowbreak carlo.m} oraz \texttt{run\_\allowbreak range\_\allowbreak limit\_\allowbreak sweep.m}) i~zestaw testów regresyjnych \texttt{tests/\allowbreak verify\_\allowbreak regression.m}. Przy jej projektowaniu konsekwentnie stosowano zasadę DRY (ang.~\textit{Don't Repeat Yourself}): wszystkie wymienione skrypty wywołują bezpośrednio produkcyjne funkcje z~\texttt{src/} i~\texttt{lib/}, a~żaden z~nich nie zawiera własnej kopii kodu przetwarzania sygnałów. Gwarantuje to, że wykresy diagnostyczne z~niniejszego rozdziału, wyniki badań z~rozdziału~\ref{sec:simulation_results} oraz testy regresyjne powstają przy użyciu dokładnie tego samego kodu. Aby było to możliwe również dla wizualizacji etapów pośrednich, struktura \texttt{radar\_data} zwracana przez \texttt{step3\_receiver} udostępnia, oprócz periodogramów wynikowych, macierze pośrednie: surową transmitancję \texttt{H\_raw}, transmitancję po usunięciu ech stałych \texttt{H\_no\_static}, transmitancję po interpolacji luki DC \texttt{H\_shifted}, dwuwymiarowe okno wagowe \texttt{W} oraz współczynnik wzmocnienia szumu \texttt{zeta\_dB}.

Zależności pomiędzy modułami oraz przepływ danych w~potoku przedstawiono na rysunku~\ref{fig:software_architecture}. Na schemacie wyróżniono trzy rodzaje bloków: fasady potoku z~katalogu \texttt{src/} (ramki niebieskie), autorskie funkcje biblioteczne z~katalogu \texttt{lib/} (bloki z~ciągłym obrysem) oraz funkcje pochodzące z~pakietów narzędziowych MATLAB-a (bloki pomarańczowe z~obrysem przerywanym), co pozwala jednoznacznie oddzielić gotowe komponenty środowiska od kodu własnego zestawionego w~tabeli~\ref{tab:toolboxes_vs_custom}.

% !TeX root = ../main.tex
% Schemat blokowy architektury oprogramowania i przeplywu danych (Rys. 3.1)
\begin{figure}[H]
\centering
\resizebox{\textwidth}{!}{%
\begin{tikzpicture}[
    font=\sffamily\footnotesize,
    >=Stealth,
    node distance=5mm and 6mm,
    % Style warstw (ramki zewnetrzne)
    ctrlframe/.style={
        draw=black!55, fill=black!4, rounded corners=3pt,
        line width=0.8pt, inner sep=6pt
    },
    stageframe/.style={
        draw=blue!55!black, fill=blue!3, rounded corners=3pt,
        line width=0.9pt, inner sep=7pt
    },
    % Style blokow wewnetrznych
    ctrlbox/.style={
        draw=black!50, fill=white, rounded corners=2pt,
        line width=0.6pt, align=center, inner sep=4pt,
        minimum height=10.5mm
    },
    libbox/.style={
        draw=blue!50!black, fill=white, rounded corners=2pt,
        line width=0.65pt, align=center, inner sep=4.5pt,
        minimum height=10mm
    },
    tbbox/.style={
        draw=orange!75!black, fill=orange!10, rounded corners=2pt,
        line width=0.75pt, dashed, align=center, inner sep=4.5pt,
        minimum height=10mm
    },
    stagehdr/.style={
        font=\sffamily\footnotesize\bfseries, text=blue!55!black,
        anchor=west, inner sep=1pt
    },
    ctrlhdr/.style={
        font=\sffamily\footnotesize\bfseries, text=black!75,
        anchor=west, inner sep=1pt
    },
    dataflow/.style={
        ->, line width=0.95pt, draw=blue!55!black
    },
    internalflow/.style={
        ->, line width=0.7pt, draw=black!70
    },
    ctrlflow/.style={
        ->, line width=0.7pt, draw=black!55, densely dashed
    },
    edgelbl/.style={
        font=\sffamily\scriptsize, fill=white, inner sep=1.5pt,
        text=black!85, align=center
    }
]

% =========================================================================
% WARSTWA 0: STEROWANIE, EKSPERYMENTY I TESTY (DRY)
% =========================================================================
\node[ctrlbox, text width=2.45cm] (init) at (0, 0) {%
    \texttt{init\_paths.m}\\[1pt]
    {\scriptsize konfiguracja ścieżek}
};
\node[ctrlbox, text width=4.0cm, right=5mm of init] (scripts) {%
    \texttt{scripts/run\_random\_demo.m}\\
    \texttt{scripts/figures/*}\\[1pt]
    {\scriptsize generacja rysunków pracy}
};
\node[ctrlbox, text width=4.1cm, right=5mm of scripts] (mc) {%
    \texttt{scripts/experiments/*}\\
    {\scriptsize \texttt{run\_detection\_monte\_carlo.m}}\\
    {\scriptsize \texttt{run\_range\_limit\_sweep.m}}
};
\node[ctrlbox, text width=3.55cm, right=5mm of mc] (tests) {%
    \texttt{tests/verify\_regression.m}\\[1pt]
    {\scriptsize \textit{Golden Master} i~testy E2E}
};

\draw[ctrlflow] (init) -- (scripts);

\path (init.north) ++(0, 5.2mm) coordinate (ctrl_top_pad);
\begin{scope}[on background layer]
    \node[ctrlframe, fit=(init)(scripts)(mc)(tests)(ctrl_top_pad),
          inner ysep=5pt, inner xsep=7pt] (ctrl_frame) {};
    \node[ctrlhdr] at ([xshift=5pt, yshift=-3.5mm]ctrl_frame.north west)
        {Warstwa inicjalizacji, eksperymentów i~testów (DRY -- bez duplikacji logiki DSP)};
\end{scope}

% =========================================================================
% ETAP 1: NADAJNIK (LEWA KOLUMNA)
% =========================================================================
\node[libbox, text width=6.1cm, anchor=north west] (gwp)
    at ([xshift=7pt, yshift=-14mm]ctrl_frame.south west) {%
    \texttt{lib/get\_wifi\_params.m}\\[1pt]
    {\scriptsize parametry PHY 802.11a/ax, auto-skalowanie ramki, maska aktywna, luki DC}
};
\node[tbbox, text width=6.1cm, below=5.5mm of gwp] (wtx) {%
    \textbf{[WLAN Toolbox]}\\[1pt]
    {\scriptsize \texttt{wlanNonHTConfig} / \texttt{wlanHESUConfig}, \texttt{wlanHEOFDMInfo},}\\
    {\scriptsize \texttt{wlanFieldIndices}, \texttt{wlanSampleRate}, \texttt{transmitTime}, \texttt{wlanWaveformGenerator}}
};

\draw[internalflow] (gwp) -- (wtx);

\path (gwp.north) ++(0, 5.2mm) coordinate (s1_hdr_pad);
\begin{scope}[on background layer]
    \node[stageframe, fit=(gwp)(wtx)(s1_hdr_pad), inner xsep=6pt, inner ysep=5pt] (s1_frame) {};
    \node[stagehdr] at ([xshift=5pt, yshift=-3.6mm]s1_frame.north west)
        {1. Nadajnik: \texttt{src/step1\_transmitter.m}};
\end{scope}

% =========================================================================
% ETAP 2: KANAL PROPAGACYJNY (PRAWA KOLUMNA)
% =========================================================================
\node[libbox, text width=3.25cm, anchor=north east] (ct)
    at ([xshift=-7pt, yshift=-14mm]ctrl_frame.south east) {%
    \texttt{lib/configure\_targets.m}\\[1pt]
    {\scriptsize opóźnienia, Doppler, RCS}\\
    {\scriptsize (\texttt{get\_tgax\_propagation\_params})}
};
\node[libbox, text width=2.85cm, left=3.5mm of ct] (gcf) {%
    \texttt{lib/get\_carrier\_frequency.m}\\[1pt]
    {\scriptsize mapowanie pasma na $f_c$}
};
\node[tbbox, text width=6.75cm, below=5mm of ct.south east, anchor=north east] (wch) {%
    \textbf{[WLAN Toolbox]} \texttt{wlanTGaxChannel} {\scriptsize (Model-A/B/C/E: DPI + clutter)}\\[1pt]
    {\scriptsize + autorska synteza celów ruchomych (opóźnienie ułamkowe FFT + Doppler)}
};
\node[tbbox, text width=6.75cm, below=4.5mm of wch] (ctb) {%
    \textbf{[Communications Toolbox]} \texttt{awgn()}\\[1pt]
    {\scriptsize szum termiczny zakotwiczony względem mocy nominalnej}
};

\draw[internalflow] (gcf.south) -- (gcf.south |- wch.north);
\draw[internalflow] (ct.south) -- (ct.south |- wch.north);
\draw[internalflow] (wch) -- (ctb);

\path (ct.north) ++(0, 5.2mm) coordinate (s2_hdr_pad);
\begin{scope}[on background layer]
    \node[stageframe, fit=(gcf)(ct)(wch)(ctb)(s2_hdr_pad), inner xsep=6pt, inner ysep=5pt] (s2_frame) {};
    \node[stagehdr] at ([xshift=5pt, yshift=-3.6mm]s2_frame.north west)
        {2. Kanał: \texttt{src/step2\_channel.m}};
\end{scope}

% =========================================================================
% ETAP 3: ODBIORNIK - PRZETWARZANIE SYGNALU (PELNA SZEROKOSC)
% =========================================================================
\node[tbbox, text width=4.8cm, anchor=north west] (wsync)
    at ([xshift=7pt, yshift=-14.5mm]s2_frame.south -| ctrl_frame.west) {%
    \textbf{[WLAN Toolbox]} \texttt{wlanPacketDetect}\\
    + \texttt{wlanSymbolTimingEstimate}\\[1pt]
    {\scriptsize + autorski \textit{leading-edge} (L-LTF) / wyzwalacz sprzętowy}
};
\node[libbox, text width=3.4cm, right=5.5mm of wsync] (demod) {%
    \texttt{lib/demodulate.m}\\[1pt]
    {\scriptsize usunięcie CP + kolumnowe FFT}\\
    {\scriptsize (tor Rx i~tor odniesienia Tx)}
};
\node[libbox, text width=5.6cm, right=5.5mm of demod] (est) {%
    \texttt{lib/estimate\_channel.m}\\[1pt]
    {\scriptsize estymacja ZF, filtr MTI ($-\mathrm{mean}$), $\zeta_{\mathrm{ZF}}$}\\
    {\scriptsize wywołuje: \texttt{lib/interpolate\_dc\_nulls.m}}
};

% Drugi rzad Etapu 3 (przeplyw w lewo: est -> rdm -> ax)
\node[libbox, text width=7.5cm, below=6mm of est.south east, anchor=north east] (rdm) {%
    \texttt{lib/compute\_range\_doppler\_map.m}\\[1pt]
    {\scriptsize \textcolor{orange!75!black}{\textbf{[Signal Processing Toolbox]}} \texttt{blackmanharris} (okno 2D)}\\
    {\scriptsize + 2D IFFT/FFT z~dopełnianiem zerami ($4\times$ ZP)}
};
\node[libbox, text width=6.4cm, anchor=west] (ax)
    at (wsync.west |- rdm.center) {%
    \texttt{lib/calculate\_physical\_axes.m}\\[1pt]
    {\scriptsize kalibracja osi fizycznych $(R, v)$ i~złożenie \texttt{radar\_data}}
};

\draw[internalflow] (wsync) -- (demod);
\draw[internalflow] (demod) -- (est);
\draw[internalflow] (est.south) -- (est.south |- rdm.north);
\draw[internalflow] (rdm) -- (ax);

\path (wsync.north) ++(0, 5.2mm) coordinate (s3_hdr_pad);
\begin{scope}[on background layer]
    \node[stageframe, fit=(wsync)(demod)(est)(rdm)(ax)(s3_hdr_pad), inner xsep=6pt, inner ysep=5pt] (s3_frame) {};
    \node[stagehdr] at ([xshift=5pt, yshift=-3.6mm]s3_frame.north west)
        {3. Odbiornik -- przetwarzanie sygnału: \texttt{src/step3\_receiver.m}};
\end{scope}

% =========================================================================
% ETAP 4: ODBIORNIK - DETEKCJA I EKSTRAKCJA (PELNA SZEROKOSC)
% =========================================================================
\node[libbox, text width=3.7cm, anchor=north west] (reg)
    at ([xshift=7pt, yshift=-17mm]s3_frame.south west) {%
    \texttt{lib/get\_detection\_algorithm.m}\\[1pt]
    {\scriptsize centralny rejestr detektorów}\\
    {\scriptsize (\texttt{'binary'} / \texttt{'coherent'})}
};
\node[libbox, text width=6.5cm, anchor=north west] (bstc)
    at ([xshift=10mm, yshift=-11mm]s3_frame.south -| reg.east) {%
    \texttt{lib/binary\_successive\_target\_cancellation.m}\\[1pt]
    {\scriptsize próg CFAR + maskowanie elipsą \textit{null-to-null}}
};
\node[libbox, text width=6.5cm, below=4.5mm of bstc] (cstc) {%
    \texttt{lib/coherent\_successive\_target\_cancellation.m}\\[1pt]
    {\scriptsize próg CFAR + Krok 4 (rezydua) + replika 2D-FFT}
};
\node[libbox, text width=3.1cm, anchor=east] (pnfr)
    at ([xshift=-7pt]s3_frame.east |- reg.center) {%
    \texttt{lib/compute\_pnfr.m}\\[1pt]
    {\scriptsize wyznaczenie kontrastu}\\
    {\scriptsize PNFR wykrytych celów}
};

\draw[internalflow] (reg.east) -- ++(2.5mm, 0) |- (bstc.west)
    node[edgelbl, pos=0.72, above=1pt] {\texttt{'binary'}};
\draw[internalflow] (reg.east) -- ++(2.5mm, 0) |- (cstc.west)
    node[edgelbl, pos=0.72, below=1pt] {\texttt{'coherent'}};
\draw[internalflow] (bstc.east) -- ++(2.5mm, 0) |- (pnfr.west);
\draw[internalflow] (cstc.east) -- ++(2.5mm, 0) |- (pnfr.west);

\path (bstc.north) ++(0, 5.2mm) coordinate (s4_hdr_pad);
\begin{scope}[on background layer]
    \node[stageframe, fit=(reg)(bstc)(cstc)(pnfr)(s4_hdr_pad), inner xsep=6pt, inner ysep=5pt] (s4_frame) {};
    \node[stagehdr] at ([xshift=5pt, yshift=-3.6mm]s4_frame.north west)
        {4. Odbiornik -- detekcja i~ekstrakcja celów: \texttt{src/step4\_interpreter.m}};
\end{scope}

% =========================================================================
% POLACZENIA MIEDZYETAPOWE (PRZEPLYW DANYCH)
% =========================================================================
% Warstwa sterujaca -> Etap 1
\draw[dataflow] ([xshift=1.8cm]ctrl_frame.south west) -- ([xshift=1.8cm]s1_frame.north west)
    node[edgelbl, midway, right=2pt] {parametry scenariusza (\texttt{standard}, \texttt{cbw}, \texttt{mcs}, \texttt{targets})};

% Etap 1 -> Etap 2
\draw[dataflow] (s1_frame.east |- gwp.south) -- (s2_frame.west |- gwp.south)
    node[edgelbl, midway, above=2pt] {\texttt{waveform}, \texttt{fs},\\ \texttt{params}, \texttt{cfg}};

% Etap 1 -> Etap 3 (tor odniesienia Tx w ukladzie wspollokowanym)
\draw[dataflow] (s1_frame.south) -- (s1_frame.south |- s3_frame.north)
    node[edgelbl, midway] {\texttt{waveform}, \texttt{fs}, \texttt{params} (sygnał odniesienia Tx)};

% Etap 2 -> Etap 3 (tor odbiorczy Rx)
\draw[dataflow] (s2_frame.south) -- (s2_frame.south |- s3_frame.north)
    node[edgelbl, midway] {\texttt{signal}, \texttt{fc}, (\texttt{intrinsic\_delay})};

% Etap 3 -> Etap 4
\draw[dataflow] (s3_frame.south) -- (s4_frame.north)
    node[edgelbl, midway] {\texttt{radar\_data} (\texttt{CPer\_F}, \texttt{Per\_F}, \texttt{W}, \texttt{H\_*}, osie $R, v$)};

% Etap 4 -> Warstwa sterujaca / wizualizacyjna (petla powrotna po prawej krawedzi)
\draw[dataflow] (s4_frame.east) -- ++(5.5mm, 0) |- (ctrl_frame.east)
    node[edgelbl, pos=0.26, rotate=90, anchor=south]
        {\texttt{targets}, \texttt{CPer\_F}, \texttt{Per\_F}, \texttt{eta}, \texttt{pnfr}, \texttt{diag\_info}};

% =========================================================================
% LEGENDA
% =========================================================================
\node[anchor=north, yshift=-4.5mm] (legend) at (s4_frame.south) {%
    \begin{tikzpicture}[baseline=-0.5ex, font=\sffamily\scriptsize]
        \node[ctrlframe, minimum width=6.5mm, minimum height=3.5mm, inner sep=0pt] (l1) at (0,0) {};
        \node[right=1.5mm of l1] (t1) {Warstwa sterująca i~testowa (DRY)};
        \node[stageframe, minimum width=6.5mm, minimum height=3.5mm, inner sep=0pt, right=5mm of t1] (l2) {};
        \node[right=1.5mm of l2] (t2) {Fasada potoku (\texttt{src/})};
        \node[libbox, minimum width=6.5mm, minimum height=3.5mm, inner sep=0pt, right=5mm of t2] (l3) {};
        \node[right=1.5mm of l3] (t3) {Kod autorski DSP (\texttt{lib/})};
        \node[tbbox, minimum width=6.5mm, minimum height=3.5mm, inner sep=0pt, right=5mm of t3] (l4) {};
        \node[right=1.5mm of l4] (t4) {Pakiet narzędziowy MATLAB (Toolbox)};
    \end{tikzpicture}
};

\end{tikzpicture}%
}
\caption{Schemat blokowy architektury oprogramowania symulatora radaru pasywnego IEEE~802.11 oraz przepływu danych między modułami autorskimi i~pakietami narzędziowymi MATLAB.}
\label{fig:software_architecture}
\end{figure}


Przetwarzanie przebiega sekwencyjnie w~czterech etapach. Nadajnik na podstawie parametrów scenariusza (standard, szerokość kanału, indeks MCS) generuje zespolony przebieg pasma podstawowego \texttt{waveform} wraz z~częstotliwością próbkowania \texttt{fs}, strukturą parametrów warstwy fizycznej \texttt{params} i~obiektem konfiguracyjnym WLAN Toolbox \texttt{cfg}. Moduł kanału przyjmuje przebieg, profil kanału TGax oraz opis celów, a~zwraca sygnał odebrany \texttt{signal}, częstotliwość nośną \texttt{fc} oraz opóźnienie własne filtru kanału \texttt{intrinsic\_delay}. Część przetwarzająca odbiornika przekształca sygnał odebrany w~strukturę \texttt{radar\_data}, zawierającą zespolony periodogram \texttt{CPer\_F}, periodogram mocy \texttt{Per\_F}, okno \texttt{W}, macierze pośrednie oraz skalibrowane osie odległości i~prędkości. Ostatni etap zwraca do warstwy sterującej listę wykrytych celów \texttt{targets} wraz z~estymatami odległości, prędkości i~skutecznej powierzchni odbicia, periodogramy rezydualne po ekstrakcji celów, próg detekcji \texttt{eta}, metryki kontrastu \texttt{pnfr} oraz strukturę diagnostyczną \texttt{diag\_info}. Struktura \texttt{params} jest przekazywana z~nadajnika bezpośrednio do odbiornika, dzięki czemu odbiornik korzysta z~identycznej maski aktywnych podnośnych, indeksów luk DC i~położenia pola danych, jakie wyznaczono przy generacji ramki. Dla zachowania możliwości samodzielnego użycia funkcja \texttt{step3\_receiver} akceptuje alternatywnie samą nazwę standardu i~odtwarza wówczas parametry za pomocą \texttt{get\_wifi\_params}.

Układ połączeń na rysunku~\ref{fig:software_architecture} wynika bezpośrednio z~geometrii systemu opisanej w~podrozdziale~\ref{subsec:pbr_fundamentals}. W~rozważanym układzie współlokowanym, z~dwiema odizolowanymi antenami nadawczą i~odbiorczą, odbiornik znajduje się przy nadajniku i~ma dostęp do transmitowanego przebiegu. Z~tego względu \texttt{step3\_receiver} otrzymuje dwa niezależne wejścia: bezszumowy sygnał odniesienia \texttt{waveform} bezpośrednio z~\texttt{step1\_transmitter} oraz sygnał odebrany \texttt{signal} z~\texttt{step2\_channel}. Oba tory są demodulowane tą samą funkcją \texttt{demodulate}, a~ich ilorazem jest estymowana transmitancja kanału. Wspólne źródło zegara nadajnika i~odbiornika pozwala ponadto wyznaczyć początek ramki w~sygnale odebranym na podstawie znanego przesunięcia toru (argument \texttt{sync\_idx}, w~symulacji równy \texttt{intrinsic\_delay}), co odpowiada wyzwalaczowi sprzętowemu w~rzeczywistym urządzeniu współlokowanym.

Nawigacyjne zestawienie wszystkich plików potoku wraz z~ich odpowiedzialnością, głównymi argumentami oraz odniesieniami do teorii przedstawiono w~tabeli~\ref{tab:software_modules}. Szczegółowy opis implementacji poszczególnych bloków zawierają podrozdziały~\ref{subsec:transmitter_implementation}--\ref{subsec:receiver_implementation}.

\begin{table}[H]
    \centering
    \linespread{1.05}\selectfont
    \caption{Moduły programowe symulatora wraz z~ich odpowiedzialnością, interfejsem i~odniesieniem do teorii}
    \label{tab:software_modules}
    \scriptsize
    \begin{tabular}{|p{3.3cm}|p{4.5cm}|p{4.3cm}|p{2.2cm}|}
        \hline
        \textbf{Plik} & \textbf{Główna odpowiedzialność} & \textbf{Kluczowe wejścia $\rightarrow$ wyjścia} & \textbf{Teoria} \\
        \hline
        \multicolumn{4}{|l|}{\textit{Fasady potoku (\texttt{src/})}} \\
        \hline
        \texttt{step1\_\allowbreak transmitter.m} & Synteza ramki PPDU IEEE~802.11a/ax z~losowym polem danych & \texttt{WIFI\_STANDARD}, \texttt{cbw}, \texttt{mcs}, \texttt{psdu\_length} $\rightarrow$ \texttt{waveform}, \texttt{fs}, \texttt{params}, \texttt{cfg} & \ref{subsec:wifi_signal}, tab.~\ref{tab:wifi_params} \\
        \hline
        \texttt{step2\_\allowbreak channel.m} & Kanał TGax (DPI i~echa stałe), synteza ech celów ruchomych, szum termiczny wg budżetu mocy & \texttt{waveform}, \texttt{fs}, \texttt{model}, \texttt{cfg}, \texttt{targets}, \texttt{p\_rx\_dBm}, \texttt{nf\_dB} $\rightarrow$ \texttt{signal}, \texttt{fc}, \texttt{intrinsic\_delay} & \ref{subsec:pbr_fundamentals} \\
        \hline
        \texttt{step3\_\allowbreak receiver.m} & Synchronizacja ramki, demodulacja, estymacja i~kondycjonowanie transmitancji, mapa odległość--prędkość & \texttt{signal}, \texttt{waveform}, \texttt{fs}, \texttt{fc}, \texttt{params}, \texttt{sync\_idx} $\rightarrow$ \texttt{radar\_data} & \ref{subsec:ofdm_modulation}, \ref{subsec:channel_estimation}--\ref{subsec:range_doppler_map} \\
        \hline
        \texttt{step4\_\allowbreak interpreter.m} & Wybór detektora, progowanie CFAR, ekstrakcja celów, metryka PNFR & \texttt{radar\_data}, \texttt{p\_FA}, \texttt{stc\_method} $\rightarrow$ \texttt{targets}, \texttt{CPer\_F}, \texttt{Per\_F}, \texttt{eta}, \texttt{pnfr}, \texttt{diag\_info} & \ref{subsec:target_detection_extraction} \\
        \hline
        \multicolumn{4}{|l|}{\textit{Parametryzacja nadajnika i~kanału (\texttt{lib/})}} \\
        \hline
        \texttt{get\_wifi\_\allowbreak params.m} & Obiekt \texttt{cfg}, $N_{\mathrm{FFT}}$, $N_{\mathrm{CP}}$, automatyczne skalowanie długości ładunku, maska aktywnych podnośnych, indeksy luk DC & \texttt{standard}, \texttt{psdu\_length}, \texttt{cbw}, \texttt{mcs} $\rightarrow$ \texttt{params} & \ref{subsec:wifi_signal}, tab.~\ref{tab:wifi_params} \\
        \hline
        \texttt{get\_carrier\_\allowbreak frequency.m} & Odwzorowanie nazwy pasma na częstotliwość nośną & \texttt{band\_input} $\rightarrow$ \texttt{fc}, \texttt{band\_str} & \ref{subsec:pbr_fundamentals} \\
        \hline
        \texttt{configure\_\allowbreak targets.m} & Przeliczenie opisu fizycznego celu na opóźnienie, przesunięcie Dopplera i~amplitudę echa & \texttt{targets}, \texttt{fs}, \texttt{fc}, \texttt{channel\_model}, \texttt{p\_rx\_dBm} $\rightarrow$ \texttt{delays}, \texttt{dopplers}, \texttt{taps} & \ref{subsec:pbr_fundamentals} \\
        \hline
        \texttt{get\_tgax\_\allowbreak propagation\_\allowbreak params.m} & Parametry dwuodcinkowego modelu tłumienia TGn/TGax (odległość przełamania, wykładniki LOS/NLOS) & \texttt{model} $\rightarrow$ \texttt{p} (\texttt{d\_bp}, \texttt{n\_los}, \texttt{n\_nlos}) & \ref{subsec:channel_implementation} \\
        \hline
        \multicolumn{4}{|l|}{\textit{Przetwarzanie sygnału w~odbiorniku (\texttt{lib/})}} \\
        \hline
        \texttt{demodulate.m} & Usunięcie prefiksu cyklicznego i~kolumnowa transformata FFT & \texttt{signal}, \texttt{n\_fft}, \texttt{n\_cp} $\rightarrow$ \texttt{F\_matrix}, \texttt{n\_syms} & \ref{subsec:ofdm_modulation} \\
        \hline
        \texttt{estimate\_\allowbreak channel.m} & Estymacja ZF, współczynnik $\zeta_{\mathrm{ZF}}$, filtr MTI w~czasie wolnym, wywołanie interpolacji luki DC & \texttt{F\_rx}, \texttt{F\_tx}, \texttt{params} $\rightarrow$ \texttt{H\_shifted}, \texttt{H\_raw}, \texttt{zeta\_dB}, \texttt{H\_no\_static} & \ref{subsec:channel_estimation}, \ref{subsec:clutter_filtering} \\
        \hline
        \texttt{interpolate\_\allowbreak dc\_nulls.m} & Liniowa interpolacja spójnych luk zerowych podnośnych wewnętrznych & \texttt{H}, \texttt{dc\_indices} $\rightarrow$ \texttt{H} & \ref{subsec:clutter_filtering} \\
        \hline
        \texttt{compute\_range\_\allowbreak doppler\_map.m} & Okno 2D Blackmana--Harrisa, dopełnianie zerami, periodogram zespolony i~periodogram mocy & \texttt{H\_shifted}, \texttt{Nfft}, \texttt{n\_symbols}, \texttt{active\_mask} $\rightarrow$ \texttt{CPer\_F}, \texttt{W}, \texttt{Per\_F}, \texttt{N\_per}, \texttt{M\_per} & \ref{subsec:range_doppler_map} \\
        \hline
        \texttt{calculate\_\allowbreak physical\_axes.m} & Odwzorowanie indeksów DFT/IDFT na osie odległości i~prędkości & \texttt{fs}, \texttt{fc}, \texttt{Nfft}, \texttt{Ncp}, \texttt{N\_per}, \texttt{M\_per} $\rightarrow$ \texttt{axis\_range}, \texttt{axis\_velocity} & \ref{subsec:pbr_fundamentals}, \ref{subsec:range_doppler_map} \\
        \hline
        \multicolumn{4}{|l|}{\textit{Detekcja i~ocena wyników (\texttt{lib/})}} \\
        \hline
        \texttt{get\_detection\_\allowbreak algorithm.m} & Centralny rejestr detektorów (identyfikator, nazwa, uchwyt wykonawczy) & \texttt{method} $\rightarrow$ \texttt{alg} & \ref{subsec:target_detection_extraction} \\
        \hline
        \texttt{binary\_\allowbreak successive\_\allowbreak target\_\allowbreak cancellation.m} & Próg CFAR, iteracyjna detekcja z~maskowaniem binarnym listka głównego & \texttt{Per\_F}, \texttt{W}, \texttt{p\_FA} $\rightarrow$ \texttt{targets}, \texttt{B}, \texttt{eta}, \texttt{diag\_info} & \ref{subsec:target_detection_extraction} \\
        \hline
        \texttt{coherent\_\allowbreak successive\_\allowbreak target\_\allowbreak cancellation.m} & Próg CFAR, iteracyjna detekcja z~koherentnym odejmowaniem repliki celu & \texttt{CPer\_F}, \texttt{Per\_F}, \texttt{W}, \texttt{p\_FA} $\rightarrow$ \texttt{targets}, \texttt{CPer\_F}, \texttt{eta}, \texttt{diag\_info} & \ref{subsec:target_detection_extraction} \\
        \hline
        \texttt{compute\_pnfr.m} & Kontrast szczytu celu względem tła szumowo-zakłóceniowego (PNFR) & \texttt{CPer}, \texttt{target\_idx}, \texttt{guard} $\rightarrow$ \texttt{metrics} & \ref{subsec:target_detection_extraction} \\
        \hline
    \end{tabular}
\end{table}

\subsection{Implementacja nadajnika sygnału IEEE 802.11}
\label{subsec:transmitter_implementation}

\subsection{Implementacja modelu kanału propagacyjnego i~celów ruchomych}
\label{subsec:channel_implementation}

\subsection{Implementacja odbiornika radarowego}
\label{subsec:receiver_implementation}
